\section{Related Work}
\label{sec:related}

\paragraph{Automated scientific discovery and question generation.}
Computational discovery has a long lineage, from rule-based rediscovery
of physical laws \citep{langley1987} through closed-loop robot
scientists \citep{king2009} to the Nobel Turing Challenge's call for
AI scientists \citep{kitano2021}. Recent LLM-based systems generate
research ideas, hypotheses, or full papers: literature-based generation
\citep{wang2024scimon}, agentic idea refinement
\citep{baek2024researchagent}, and end-to-end automated research
\citep{lu2024aiscientist}. Literature-based discovery pioneered the
underlying intuition that recombining published evidence can anticipate
findings later verified empirically \citep{swanson1986}. Our work is
orthogonal to all of these: we do not propose a better generator; we
propose the missing evaluation.

\paragraph{Evaluating generated ideas.}
Existing evaluations are dominated by human preference and LLM scoring.
\citet{si2024llmideas} ran a large expert study comparing human and LLM
research ideas on rated novelty and excitement---the most rigorous
instance of the expert-score paradigm, and still a measurement of
\emph{opinion at generation time} rather than of what the ideas turned
out to be worth. LLM-as-judge scoring inherits known biases (position,
verbosity, self-preference) and, for questions about the future,
cannot be validated against ground truth at all. Historical backtesting
replaces both with an outcome variable that exists independently of any
rater: the subsequent behavior of the scientific community.

\paragraph{Backtesting and forecasting benchmarks.}
Scoring a strategy on held-out history is standard in quantitative
finance, along with well-documented failure modes---overfitting to the
backtest itself \citep{bailey2014backtest}---that motivate our
freezing and no-overwrite rules. Forecasting benchmarks score models on
events that resolve after training \citep{zou2022autocast}; retrodictive
evaluation with temporal holdouts is likewise used to test whether
models anticipate later discoveries \citep{swanson1986,king2009}.
We transplant this design to a harder target: not whether a stated
event occurs, but whether an open-ended research question earns the
community's future investment. Code benchmarks such as SWE-bench
\citep{jimenez2024swebench} demonstrated how a well-specified task
format plus frozen data can reorganize a research area around
measurable progress; we aim the same mechanism at question discovery.

\paragraph{Data contamination.}
Temporal splits are increasingly used to control LLM memorization in
evaluation. Our protocol controls the \emph{retrieval} channel
completely (corpus manifests, isolation rules, CI checks) and treats the
\emph{weights} channel---models whose training data postdates the
cutoff---as a declared, audited threat rather than a solved problem
(Section~\ref{sec:threats}).
